\clearpage
\section{Recursive copying residuals and genuine specimens}
\label{sec:copying-residual-specimens}
The copying diagrams of Section~\ref{sec:copying-box-permutation-overlaps}
now have genuine recursive filling certificates. Both surviving branches are
represented by the original \texttt{DeepFillingEdge} constructors, and the
resulting closed residuals give the original specimen system's quasi-specimens.
There is also a forward and converse robustness-case theorem for closed
all-deep copying, with arbitrary certified bodies and every genuine unary
contraction-tree presentation. The structural certificate construction itself
allows arbitrary tree arity, arbitrary free variable and atom interfaces, and
arbitrary nesting depth.

The mathematical advance is the passage from a structural derivation to the
particular recursive certificate consumed by the robustness framework.
Previously it was not enough to know that a residual had been assembled from
BPerm pairs: the required common context, paired filling families and token
coordinates still had to be supplied. They are now constructed explicitly.
The assembly and transfer results admit simultaneous shallow and deep
contributions. They do not prove that every overlap of an arbitrary original
specimen admits the required decomposition. The local deep beta
case is completed in Section~\ref{sec:copying-beta-residuals}; the remaining
operation transfers, exhaustive decomposition and universal forward/converse
BPerm robustness assembly remain unfinished. Proposition~12.2 is
therefore still unproved. No program, translation clause, stability premise,
copyability premise or machine rule is changed.

\subsection{Flattening nested expansion}
\termdef{recursive-expansion-flattening}
Let $C$ be a typed finite host, $F_e$ a graph at each host edge, and $H_{e,f}$
a graph at each edge of $F_e$. Write $C[F]$ for simultaneous expansion. The
flattened family is $\widehat H_{(e,f)}=H_{e,f}$. There is an ordered,
boundary-fixing isomorphism
\[
 (C[F])[\widehat H]\cong C[e\mapsto F_e[H_e]].
\]
Its internal vertices are reassociated by
\[
 (V_C+\textstyle\sum_e V_{F_e})+
   \textstyle\sum_{(e,f)}V_{H_{e,f}}
 \cong V_C+\textstyle\sum_e(V_{F_e}+\sum_f V_{H_{e,f}}),
\]
and its edge bijection sends $((e,f),g)$ to $(e,(f,g))$.
The definitions \texttt{flat}, \texttt{vertices} and \texttt{edges} give these
maps and their inverses. Their two inverse laws are immediate on each tagged
summand. All external positions remain fixed.

\lean{Graph/ExpansionAssociativity}{14}{14}{ExpansionAssoc.flat}
\lean{Graph/ExpansionAssociativity}{16}{27}{ExpansionAssoc.vertices}
\lean{Graph/ExpansionAssociativity}{29}{34}{ExpansionAssoc.edges}

\Needspace{12\baselineskip}
The host-coordinate equation says that embedding a vertex of $C$ twice agrees
with embedding it once after reassociation. The middle-coordinate equation
says the same for a vertex of $F_e$. For internal vertices these are the
corresponding tag equations. At inputs and outputs, the middle equation reduces
to the host equation at the relevant source or target port. The two port lemmas
unfold incidence-list mapping and identify the port of $(e,f)$ with the insertion
of the corresponding port of $f$.

\lean{Graph/ExpansionAssociativity}{36}{39}{ExpansionAssoc.map_host}
\lean{Graph/ExpansionAssociativity}{41}{48}{ExpansionAssoc.map_middle}
\lean{Graph/ExpansionAssociativity}{50}{52}{ExpansionAssoc.sourcePort}
\lean{Graph/ExpansionAssociativity}{54}{56}{ExpansionAssoc.targetPort}

The insertion equation now handles a vertex of $H_{e,f}$: internal vertices
follow reassociation, whereas boundary vertices use a port lemma followed by
the middle equation. Finally \texttt{iso} combines the two carrier bijections.
Vertex and edge labels are unchanged; mapping each ordered source or target
list gives equality pointwise by the insertion equation. This proof imposes
no acyclicity or absence-of-self-loops condition on the host.

\lean{Graph/ExpansionAssociativity}{59}{69}{ExpansionAssoc.map_insert}
\lean{Graph/ExpansionAssociativity}{72}{87}{ExpansionAssoc.iso}

\subsection{Expansion through finite parallel assembly}
\termdef{parallel-filling-assembly}
Let $C_a$ be a finite family of contexts with a chosen ordered boundary shuffle
$B$, and let $F_{a,e}$ be their edge fillings. Expansion commutes with their
original parallel assembly:
\[
 (\mathop{\|}\nolimits_B C_a)[(a,e)\mapsto F_{a,e}]
 \cong \mathop{\|}\nolimits_B(C_a[F_a]).
\]
The definition \texttt{family} supplies the flattened family. The vertex
bijection sends a tagged host vertex to the same component's host summand,
and a tagged inserted vertex to that component's inserted summand. The edge
bijection reassociates $((a,e),g)$ to $(a,(e,g))$. The explicit inverse maps
verify both inverse equations by cases.

\lean{Graph/ParallelExpansion}{14}{14}{ParallelExpansion.family}
\lean{Graph/ParallelExpansion}{16}{21}{ParallelExpansion.vertices}
\lean{Graph/ParallelExpansion}{23}{28}{ParallelExpansion.edges}

The host embedding law respects the original shuffle at inputs and outputs.
The source- and target-port equations unfold the mapped incidence lists.
For an inserted boundary vertex these equations reduce compatibility to the
host embedding law; inserted internal vertices are immediate. The final
isomorphism checks ordered incidence by mapping the insertion law over each
list. Neither the shuffle nor the enumeration of components is replaced.

\lean{Graph/ParallelExpansion}{30}{34}{ParallelExpansion.map_host}
\lean{Graph/ParallelExpansion}{36}{38}{ParallelExpansion.sourcePort}
\lean{Graph/ParallelExpansion}{40}{42}{ParallelExpansion.targetPort}
\lean{Graph/ParallelExpansion}{44}{55}{ParallelExpansion.map_insert}
\lean{Graph/ParallelExpansion}{57}{72}{ParallelExpansion.iso}

\subsection{Explicit recursive filling witnesses}
\termdef{recursive-filling-witness}
For a pre-template $T$ and graphs $G,H:I\to O$, a
\texttt{FillingWitness} consists of a focus-free common context $C$, two
edge families $L,R$, an original \texttt{DeepFillingEdge} certificate at every
edge of $C$, and boundary-fixing recursive isomorphisms
\[
 C[L]\cong G,\qquad C[R]\cong H.
\]
The edge certificates distinguish retained corollas, actual template pairs
at holes, and recursive fillings within boxes. This structure is proof data
about graph assembly, not a new semantic relation or an extra constructor of
\texttt{DeepFilling}. The lemma \texttt{paired} packages the original expansion
constructor. The lemma \texttt{focusFree} applies its established recursive
focus-freedom theorem and transports that property along the endpoint
isomorphisms. The construction \texttt{congr} composes only those endpoint
isomorphisms, preserving all common-context data.

\lean{Copying/FillingWitness}{12}{21}{FillingWitness}
\lean{Copying/FillingWitness}{26}{28}{FillingWitness.paired}
\lean{Copying/FillingWitness}{31}{33}{FillingWitness.focusFree}
\lean{Copying/FillingWitness}{36}{38}{FillingWitness.congr}

\Needspace{12\baselineskip}
The reflexive witness retains each original edge as its own labelled corolla;
corolla collapse gives the endpoint isomorphisms. An existing
\texttt{DeepFilling} proof with a focus-free context and specified endpoint
isomorphisms also proves existence of a witness. The definition
\texttt{ofDeepFilling} selects such a witness using classical choice; its type
asserts the endpoints, not definitional equality with the supplied context.
For an actual $T$-pair, \texttt{of\_rel} is more explicit: it uses a single
hole corolla with the supplied natural-number name, fills it on the two sides,
and collapses the surrounding corolla. Thus the chosen name is real context
data. Equal-interface reindexing, in both the graph and translation APIs,
follows by equality elimination and introduces no new filling rule.

\lean{Copying/FillingWitness}{41}{48}{FillingWitness.refl}
\lean{Copying/FillingWitness}{51}{53}{FillingWitness.ofDeepFilling}
\lean{Copying/FillingWitness}{56}{73}{FillingWitness.of_rel}
\lean{Copying/FillingWitness}{75}{77}{FillingWitness.reboundary}
\lean{Copying/FillingWitness}{79}{81}{FillingWitness.interfaces}

\subsection{Assembly without new certificate axioms}
Suppose each edge $e$ of a host $C$ has a witness for $G_e,H_e$. First expand
$C$ by the witness contexts. Its edges are pairs $(e,f)$, so use the original
left and right family at $f$ and the original recursive edge certificate there.
The new context is focus-free because each inserted context is focus-free.
Flattening followed by componentwise endpoint isomorphisms proves that the
resulting endpoints are $C[G]$ and $C[H]$. The host itself may contain holes:
all its edges have been expanded by the supplied focus-free contexts.

\lean{Copying/FillingAssembly}{11}{28}{FillingWitness.expand}

Finite parallel assembly uses the same argument and the parallel-expansion
isomorphism. Tensor is its Boolean-family instance with the original ordered
concatenation. Serial assembly expands the existing two-edge serial pattern;
the two components' interfaces already provide the required typing. Consequently
one component may contain shallow pairs and another arbitrarily nested pairs.

\lean{Copying/FillingAssembly}{31}{48}{FillingWitness.parallel}
\lean{Copying/FillingAssembly}{51}{57}{FillingWitness.tensor}
\lean{Copying/FillingAssembly}{60}{66}{FillingWitness.serial}

For a box, keep the original bound-variable, free-variable and atom counts,
and place the witness context inside that box. Its sole edge gets the existing
\texttt{boxed} certificate, whose inner proof is \texttt{paired}. The left
and right labels contain the actual expanded bodies. Relabelled corolla
collapse and congruence of boxing identify them with the intended endpoint
boxes. The construction does not evaluate inside a box or erase its body.

\lean{Copying/FillingAssembly}{69}{82}{FillingWitness.boxed}

\subsection{Both branches of the literal copying rule}
\termdef{copying-residual-filling}
Let $T$ have $n+1$ term inputs. A certified pair $G,H:*\to E_{k,h}$ has a
certified source pair by serially preceding it with the unchanged contraction
tree. The tree's primitive-label theorem supplies its focus-free reflexive
witness. To certify the two-copy assembly, independently supply an active and
a retained witness for the same endpoint pair. Their contexts and hole names
may differ. Parallel assembly combines the active witness with the original
residual tree followed by the retained witness, then applies the original
boundary reindexing. The residual tree's focus-freedom follows from its
contraction and weakening labels.

\lean{Copying/CopyFillingWitness}{11}{15}{FillingWitness.copyingBefore}
\lean{Copying/CopyFillingWitness}{19}{27}{FillingWitness.copies}

The complete \texttt{copyingAfter} witness appends the unchanged two-use term
and atom distributors. Their label invariants supply the reflexive witness;
both copies and both distributor sorts remain actual graph components.
The lambda-payload witness prefixes the original lambda operation to a boxed
witness, with the same interface cast as the existing translation. Applying
this construction to a body witness and using it in both copy branches gives
\texttt{deepCopy}. There is no termination assumption on either body.

\lean{Copying/CopyFillingWitness}{31}{37}{FillingWitness.copyingAfter}
\lean{Copying/CopyFillingWitness}{41}{44}{FillingWitness.lambdaPayload}
\lean{Copying/CopyFillingWitness}{47}{51}{FillingWitness.deepCopy}

For the shallow BPerm case, let $f:B\to C$ be an ordered body boundary
isomorphism with precisely the free action $p$ and fixed bound argument, and
assume $B$ focus-free. The earlier lambda-pair theorem gives a genuine
BPerm pair for the box and its externally compensated partner.
\texttt{payloadWitness} places that pair in a named hole and prefixes the
unchanged lambda head. Its focus-freedom follows from the body and interface
cast. The source witness adds the actual contraction tree and moves the
compensation to the outer boundary using the proved serial isomorphism.

\lean{Copying/CopyFillingWitness}{62}{67}{BoxPermutation.Copying.payloadWitness}
\lean{Copying/CopyFillingWitness}{70}{74}{BoxPermutation.Copying.sourceWitness}

\Needspace{10\baselineskip}
The residual witness uses names $0$ and $1$ for the two independently constructed
branch witnesses. The proved naturality of the complete copying reduct moves
the compensation to its external outputs. This supplies actual recursive
certificates for the previously proved shallow copying diagram. It does not
assume semantic soundness of the old structural derivation.
For all-deep bodies, \texttt{deep\_diamond} pairs two original copying steps
with the newly constructed residual witness. Lambda shape proves both
copyability premises; arbitrary body certificates are duplicated literally.

\lean{Copying/CopyFillingWitness}{78}{83}{BoxPermutation.Copying.residualWitness}
\lean{Copying/CopyFillingWitness}{87}{93}{BoxPermutation.Copying.deep_diamond}

\subsection{From a witness to an original quasi-specimen}
\termdef{unrestricted-bperm-contexts}
The paper proves BPerm robustness over unrestricted contexts in
\href{https://arxiv.org/html/1907.01257v5\#A8.SS3.SSS6}{Appendix~H.3.6}.
Binding-free contexts are a later requirement for certain operational templates,
not a restriction to impose on these copying residuals. Accordingly
\texttt{allContexts} is the predicate constantly true. This is an allowed-context
parameter of the existing specimen system; it does not change the machine,
its rootedness predicate, or its observation definition.

\lean{Copying/FillingSpecimens}{11}{11}{Specimens.allContexts}

For a token $q=(m,v)$ in the common context, first embed $v$ in each expansion,
then map it by the corresponding endpoint isomorphism. These are
\texttt{leftPosition} and \texttt{rightPosition}; the mode is unchanged.
Boundary-fixing isomorphisms and host insertion imply that an input position
remains that same input on both sides, by definitional equality.

\lean{Copying/FillingSpecimens}{17}{18}{FillingWitness.leftPosition}
\lean{Copying/FillingSpecimens}{19}{20}{FillingWitness.rightPosition}
\lean{Copying/FillingSpecimens}{22}{23}{FillingWitness.leftPosition_input}
\lean{Copying/FillingSpecimens}{25}{26}{FillingWitness.rightPosition_input}

For closed-interface endpoints $G,H:*\to\varnothing$, assume both are
hole-free and that the context token has term type. The witness then constructs
an original \texttt{Specimen}: its context and edge families are exactly those
already supplied; global hole-freedom of the expansions follows by the endpoint
isomorphisms. The two state comparison lemmas identify its projections with the
actual endpoint states in the representation-class machine. Their proofs use
quotient soundness for recursive state isomorphism, preserving both mode and
vertex. There is no change to the quotient relation.

\lean{Copying/FillingSpecimens}{29}{42}{FillingWitness.specimen}
\lean{Copying/FillingSpecimens}{44}{48}{FillingWitness.specimen_left}
\lean{Copying/FillingSpecimens}{50}{54}{FillingWitness.specimen_right}

\Needspace{17\baselineskip}
If the mode is not rewrite and $T$ is output-closed, these identities give a
quasi-specimen up to equality. In detail, assume the two rootedness hypotheses
from the existing definition of \texttt{Quasi}, use the constructed specimen,
and transport rootedness through the state identities. The two endpoint
relations are those same identities in the required orientations. The opposite
system uses the same specimen with exchanged projections and rootedness
proofs; no symmetry premise on $T$ is required. Specialising the token to
search at the sole input gives \texttt{initial\_quasi}.

\lean{Copying/FillingSpecimens}{58}{72}{FillingWitness.quasi}
\lean{Copying/FillingSpecimens}{76}{90}{FillingWitness.quasi_opposite}
\lean{Copying/FillingSpecimens}{93}{97}{FillingWitness.initial_quasi}

For a witness built by host expansion, a token at an original host vertex maps
exactly to that vertex's embedding in the corresponding endpoint expansion.
This follows from the flattening host equation and the componentwise
isomorphism, or directly by cases on the vertex. Thus
\texttt{expanded\_quasi} applies at every typed non-rewrite host token, including
internal host vertices. It requires hole-freedom of both whole expanded graphs,
not a claim that a local return is globally final.

\lean{Copying/FillingSpecimens}{100}{105}{FillingWitness.expand_leftPosition}
\lean{Copying/FillingSpecimens}{107}{112}{FillingWitness.expand_rightPosition}
\lean{Copying/FillingSpecimens}{116}{128}{FillingWitness.expanded_quasi}

\subsection{Framed transfer and exact one-step robustness cases}
\termdef{framed-residual-transfer}
Let $C:*\to\varnothing$ be a common host, and suppose its selected edge is
filled by a genuine rewrite on each side. Other component fillings are
unchanged within each side, but may differ between sides. Suppose the two
local targets return search to the same selected input position, and every
pair of target components has a recursive filling witness. Then
\texttt{framed\_rewrite} proves an original full-machine step on each expanded
host and an equality-up-to quasi-specimen for the residuals. Each step is an
instance of the existing simultaneous-expansion rewrite constructor. Its
return vertex is the host source port at that selected input; the port-label
law supplies its term typing. The residual follows from
\texttt{expanded\_quasi}. Components outside the rewritten slot can contain
arbitrary shallow and deep certified contributions.

\lean{Copying/FramedFillingTransfer}{14}{33}{FillingWitness.framed_rewrite}

Equality-up-to implies refinement-up-to for any cost preorders $P,R$:
substitute the equal states, then use reflexivity of machine refinement.
The lemma \texttt{quasi\_exchange} reverses equality-up-to by keeping the same
specimen and swapping its projections, roots and equality orientations.
Both facts are independent of a particular graph template.

\lean{Copying/FramedFillingTransfer}{37}{44}{FillingWitness.quasi_refines}
\lean{Copying/FramedFillingTransfer}{48}{52}{FillingWitness.quasi_exchange}

\termdef{matched-copying-robustness-case}
Now let $s\to u$ and $t\to v$ be actual representation-class machine steps,
with an equality-up-to residual between $u$ and $v$. For any independently
presented step $s\to u'$, the proved determinism theorem gives $u'=u$.
Choose $k=0$ additional steps on the left and $l=1$ on the right. Then
$Q(1+k,l)$ is $Q(1,1)$ and holds by reflexivity of the cost preorder. The
right path is the supplied step, and the residual is converted to
refinement-up-to. This proves exactly the existential, non-stuck branch of
the robustness conclusion for this matched case, with no determinism axiom.

\lean{Copying/FramedFillingTransfer}{57}{65}{FillingWitness.matching_robust_case}

\Needspace{14\baselineskip}
For the closed all-deep copying instance, take any certified bodies
$B,C:*\to E_{1,0}$ and any genuine tree $T$ with one input. The definitions
\texttt{source} and \texttt{target} are the original pre-copy graph in rewrite
mode and the original post-copy graph in search mode, each at input zero.
The \texttt{step} theorem projects the literal local copy step to the proved
class machine. The auxiliary hole-freedom lemma follows through lambda-head
assembly, the original interface cast and boxing.

\lean{Copying/DeepCopyRobustness}{11}{14}{BoxPermutation.DeepCopy.payload_closed}
\lean{Copying/DeepCopyRobustness}{19}{20}{BoxPermutation.DeepCopy.source}
\lean{Copying/DeepCopyRobustness}{21}{22}{BoxPermutation.DeepCopy.target}
\lean{Copying/DeepCopyRobustness}{25}{27}{BoxPermutation.DeepCopy.step}

Hole-freedom of each body gives hole-freedom of the copy reduct. Apply the
new deep-copy witness and the initial quasi-specimen construction to obtain
\texttt{quasi}; exchange the specimen projections for
\texttt{quasi\_opposite}. The two robustness-case theorems combine these
residuals with the two literal copying steps and the preceding matched-case
argument. They quantify over every cost-preorder triple and every independent
successor of the appropriate source. The closed-interface specialisation uses
one tree input; the underlying residual certificate and local deep diagram
retain arbitrary arity and open environments.

\lean{Copying/DeepCopyRobustness}{30}{35}{BoxPermutation.DeepCopy.quasi}
\lean{Copying/DeepCopyRobustness}{37}{42}{BoxPermutation.DeepCopy.quasi_opposite}
\lean{Copying/DeepCopyRobustness}{46}{52}{BoxPermutation.DeepCopy.robust_case}
\lean{Copying/DeepCopyRobustness}{55}{61}{BoxPermutation.DeepCopy.robust_case_opposite}

\subsection{Nontrivial and host-internal checks}
The shallow regression constructs a residual witness for the earlier nested
body with a swap of both free variables and free atoms, for every tree arity
and presentation. A second swap exchanges two free term ports with no atoms.
The seed is the unchanged translation of variable zero in a three-variable
environment; \texttt{changedSeed} applies the bound-fixing body action.
The named leaf witness is the shallow BPerm payload certificate for this pair.

\lean{Copying/CopyResidualExamples}{11}{13}{BoxPermutation.ResidualExamples.shallow}
\lean{Copying/CopyResidualExamples}{15}{16}{BoxPermutation.ResidualExamples.swap}
\lean{Copying/CopyResidualExamples}{18}{18}{BoxPermutation.ResidualExamples.seed}
\lean{Copying/CopyResidualExamples}{19}{19}{BoxPermutation.ResidualExamples.changedSeed}
\lean{Copying/CopyResidualExamples}{21}{23}{BoxPermutation.ResidualExamples.leaf}

\Needspace{20\baselineskip}
The graph \texttt{shallowBody} is a translated lambda. Applying the same
nontrivial free action gives \texttt{changedShallow}; its payload witness
is shallow at depth two. In the mixed assembly, the left component is this shallow witness and the
right is the earlier leaf witness below one additional lambda box. The graphs
\texttt{mixedLeft} and \texttt{mixedRight} retain a real binary-addition head
and the original tensor assembly. The definition \texttt{mixed} certifies both
components together. This is a structural copying regression; addition is
not being asserted to successfully consume lambda values.

\lean{Copying/CopyResidualExamples}{25}{26}{BoxPermutation.ResidualExamples.shallowBody}
\lean{Copying/CopyResidualExamples}{27}{27}{BoxPermutation.ResidualExamples.changedShallow}
\lean{Copying/CopyResidualExamples}{28}{30}{BoxPermutation.ResidualExamples.shallowPayload}
\lean{Copying/CopyResidualExamples}{32}{34}{BoxPermutation.ResidualExamples.mixedLeft}
\lean{Copying/CopyResidualExamples}{35}{38}{BoxPermutation.ResidualExamples.mixedRight}
\lean{Copying/CopyResidualExamples}{42}{43}{BoxPermutation.ResidualExamples.mixed}

The two additional original lambda payloads bind the first two of the three
free term ports, producing \texttt{closedLeftBody} and
\texttt{closedRightBody} with one remaining bound-argument port. The witness
\texttt{closedBody} follows by two lambda-payload assemblies. Hole-freedom of
the seed follows from translation; hole-freedom on its compensated side follows
from the already proved primitive-label invariance of the structural BPerm
pair. The shallow pair uses the same argument. Tensor, serial and lambda
assembly prove \texttt{mixed\_closed} and \texttt{closed\_body\_closed}.
The theorem \texttt{closed\_copy\_quasi} then checks both actual quasi-specimen
orientations after the final outer copying step. Mixed shallow/deep occurrence
inside the body is thus retained in both copies; this test does not replace
the missing exhaustive mixed-overlap theorem.

\lean{Copying/CopyResidualExamples}{45}{45}{BoxPermutation.ResidualExamples.closedLeftBody}
\lean{Copying/CopyResidualExamples}{46}{46}{BoxPermutation.ResidualExamples.closedRightBody}
\lean{Copying/CopyResidualExamples}{49}{50}{BoxPermutation.ResidualExamples.closedBody}
\lean{Copying/CopyResidualExamples}{52}{55}{BoxPermutation.ResidualExamples.seed_closed}
\lean{Copying/CopyResidualExamples}{57}{60}{BoxPermutation.ResidualExamples.shallow_closed}
\lean{Copying/CopyResidualExamples}{62}{68}{BoxPermutation.ResidualExamples.mixed_closed}
\lean{Copying/CopyResidualExamples}{71}{73}{BoxPermutation.ResidualExamples.closed_body_closed}
\lean{Copying/CopyResidualExamples}{76}{84}{BoxPermutation.ResidualExamples.closed_copy_quasi}

Finally, \texttt{host} is unary negation followed by a named closed hole at
depth five; \texttt{selected} is that hole. The family retains negation and
inserts a supplied closed graph in the selected slot. Its witness family
retains the former and uses the supplied certificate for the latter. The
\texttt{others} lemma proves equality of every unselected component by the
Boolean edge tag and uniqueness of the selected corolla edge. Expansion
preserves hole-freedom by cases on that same tag. The selected source port is
an internal vertex, as checked by \texttt{token\_internal}.

\lean{Copying/CopyResidualExamples}{86}{88}{BoxPermutation.ResidualExamples.host}
\lean{Copying/CopyResidualExamples}{89}{89}{BoxPermutation.ResidualExamples.selected}
\lean{Copying/CopyResidualExamples}{90}{96}{BoxPermutation.ResidualExamples.family}
\lean{Copying/CopyResidualExamples}{98}{103}{BoxPermutation.ResidualExamples.familyWitness}
\lean{Copying/CopyResidualExamples}{105}{110}{BoxPermutation.ResidualExamples.others}
\lean{Copying/CopyResidualExamples}{112}{118}{BoxPermutation.ResidualExamples.filled_closed}
\lean{Copying/CopyResidualExamples}{120}{121}{BoxPermutation.ResidualExamples.token_internal}

\Needspace{11\baselineskip}
The hosted diagram instantiates the actual framed-rewrite theorem with the
two literal copy steps and the mixed body's deep-copy certificate. It proves
both whole-host transitions and their residual quasi-specimen at the internal
vertex. A final witness uses the unchanged translation of self-application;
constructing its copying residual needs no totality or callback-termination
assumption. None of these finite-prefix checks asserts that self-application,
or the enclosing arithmetic examples, terminates.

\lean{Copying/CopyResidualExamples}{125}{141}{BoxPermutation.ResidualExamples.hosted_diamond}
\lean{Copying/CopyResidualExamples}{144}{146}{BoxPermutation.ResidualExamples.self_application}

\subsection{Milestone boundary}
The completed obligation is genuine recursive residual certification for the
shallow and all-deep copying diagrams, including assembly with mixed
contributions and original quasi-specimen packaging. The closed all-deep case
also has the complete forward/converse robustness conclusion. The transfer
result still takes the actual matched rewrite pair and its component
certificates as hypotheses. They must be constructed for every overlap of
an arbitrary original specimen. Section~\ref{sec:copying-beta-residuals}
subsequently completes local deep opening; the other operation cases still
need their transfer arguments. These are the precise
remaining BPerm tasks in Appendix~B.1. No whole-template robustness hypothesis
has been discharged by merely renaming a structural derivation.

The full build passes 485 jobs and all 2713 axiom audits, including all 94 new
declarations. Only \texttt{propext}, \texttt{Classical.choice} and
\texttt{Quot.sound} occur. Estimated completion of the Proposition~12.2 project
remains 60--70\% of mathematical effort, with low confidence: this milestone
closes a specific proof gap but not an entire remaining template family.
